% ============================================================================
% thecompanyinc-style.tex — the ONE canonical style for The Company, Inc. legal
% documents. Every antlers `thecompanyinc-*` template AND every legal-folder doc
% \input's THIS file (resolved via TEXINPUTS, injected by antlers' mkLegalDoc),
% so a single edit here restyles every corporate PDF uniformly. No per-doc copies.
%
% Typeface: Century Schoolbook (TeX Gyre Schola) — the U.S. Supreme Court face.
% Generalized from templates/nkc-master-lease/src/style.tex. Load order in main.tex:
%   % ============================================================================
% thecompanyinc-style.tex — the ONE canonical style for The Company, Inc. legal
% documents. Every antlers `thecompanyinc-*` template AND every legal-folder doc
% \input's THIS file (resolved via TEXINPUTS, injected by antlers' mkLegalDoc),
% so a single edit here restyles every corporate PDF uniformly. No per-doc copies.
%
% Typeface: Century Schoolbook (TeX Gyre Schola) — the U.S. Supreme Court face.
% Generalized from templates/nkc-master-lease/src/style.tex. Load order in main.tex:
%   % ============================================================================
% thecompanyinc-style.tex — the ONE canonical style for The Company, Inc. legal
% documents. Every antlers `thecompanyinc-*` template AND every legal-folder doc
% \input's THIS file (resolved via TEXINPUTS, injected by antlers' mkLegalDoc),
% so a single edit here restyles every corporate PDF uniformly. No per-doc copies.
%
% Typeface: Century Schoolbook (TeX Gyre Schola) — the U.S. Supreme Court face.
% Generalized from templates/nkc-master-lease/src/style.tex. Load order in main.tex:
%   % ============================================================================
% thecompanyinc-style.tex — the ONE canonical style for The Company, Inc. legal
% documents. Every antlers `thecompanyinc-*` template AND every legal-folder doc
% \input's THIS file (resolved via TEXINPUTS, injected by antlers' mkLegalDoc),
% so a single edit here restyles every corporate PDF uniformly. No per-doc copies.
%
% Typeface: Century Schoolbook (TeX Gyre Schola) — the U.S. Supreme Court face.
% Generalized from templates/nkc-master-lease/src/style.tex. Load order in main.tex:
%   \input{thecompanyinc-style}  ->  \subimport{./src}{defined_terms}
%   ->  \subimport{./src}{variables}  ->  \subimport{./src}{content}
% Requires \documentclass{scrartcl} (KOMA-Script).
% ============================================================================

% ---- Compatible packages --------------------------------------------------
\usepackage{tocbasic}
\usepackage[english]{babel}
\usepackage{scrlayer-scrpage}
\usepackage{xcolor}
\usepackage{ulem}
\usepackage{hanging}
\usepackage{ifthen}
\usepackage[totpages]{zref}
\usepackage{tabularx}
\usepackage{array}
\usepackage{booktabs}
\usepackage{longtable}
\usepackage{xltabular}
\usepackage{enumitem}

% ---- Century Schoolbook (the typeface used by the U.S. Supreme Court) ------
\usepackage[T1]{fontenc}
\usepackage{tgschola}
\usepackage{xspace}

% Headings and title in the serif family (Century Schoolbook), not KOMA's sans.
\addtokomafont{disposition}{\rmfamily}

% ---- Document-wide constants ----------------------------------------------
% main.tex normally sets \title directly; \documentTitle is a safe fallback.
\providecommand{\documentTitle}{Document}
\newcommand{\omitText}{Intentionally Deleted}

% ---- Defined-term machinery -----------------------------------------------
% \dtregister stores a term's display text; \dtuse renders it in small caps and
% links back to its definition; \dtdef marks the definition (the anchor).
% The generated src/defined_terms.tex (\dtregister lines + \<key> uses) is
% \input by main.tex AFTER this file, so these commands exist when it loads.
\newcommand{\dtregister}[2]{\expandafter\def\csname dtdisplay@#1\endcsname{#2}}
\newcommand{\dtuse}[1]{\hyperlink{def:#1}{\textsc{\csname dtdisplay@#1\endcsname}}\xspace}
\newcommand{\dtdef}[1]{\hypertarget{def:#1}{\textsc{\csname dtdisplay@#1\endcsname}}\xspace}

% ---- Reusable helpers -----------------------------------------------------
% Centered bold non-section header (e.g. "RECITALS", "WITNESSETH").
\newcommand{\formaltext}[1]{%
	\begingroup
	\centering
	\normalfont\normalsize\bfseries #1\par
	\endgroup
	\vspace{0.5em}
}

\newcommand{\blankspace}{\ \newline \ }

% Inline fill-in blank (for the many "____" spots in these forms).
\newcommand{\fillin}[1][2.5cm]{\rule{#1}{0.4pt}}

% A signature rule (default 3in) with an optional caption line beneath it.
\newcommand{\sigline}[1][3in]{\rule{#1}{0.4pt}}
% \sigblock{caption} -> a signature rule with a caption typeset below it.
\newcommand{\sigblock}[1]{%
	\par\noindent\sigline\par\nobreak\vspace{-0.3em}\noindent #1\par
}

% Two-up signature row: a full-width rule with two captions beneath it (left/right).
% Shared by the board/shareholder consents, whose signature grid lists the directors
% or shareholders two per line above a common signature rule.
\newcommand{\sigpair}[2]{%
	\par\vspace{2em}\noindent\rule{\textwidth}{0.4pt}\par\nobreak\vspace{0.2em}%
	\noindent\makebox[0.5\textwidth][l]{#1}\makebox[0.5\textwidth][l]{#2}\par
}

% Exhibit / Schedule heading: centered "Exhibit <letter>" + underlined title.
% Used instead of \section so exhibits stay out of the article numbering and TOC.
\newcommand{\exhibithead}[2]{%
	\begingroup
	\centering
	\normalfont\large\bfseries Exhibit #1\par
	\vspace{0.3em}
	\uline{\uppercase{#2}}\par
	\endgroup
	\vspace{1em}
}
\newcommand{\schedulehead}[2]{%
	\begingroup
	\centering
	\normalfont\large\bfseries Schedule #1\par
	\vspace{0.3em}
	\uline{\uppercase{#2}}\par
	\endgroup
	\vspace{1em}
}

% ---- Title -----------------------------------------------------------------
% Completely redefine \maketitle to avoid any KOMA errors and center a bold title.
\makeatletter
\renewcommand{\maketitle}{%
	\begingroup
	\centering
	\normalfont\large\bfseries\color{black!80}%
	\@title\par
	\vspace{0.5em}
	\endgroup
}
\makeatother

% ---- Sections: centered "ARTICLE <Roman>" + underlined uppercase title ------
\makeatletter
\renewcommand{\sectionlinesformat}[4]{%
	\Ifstr{#1}{section}{%
		\begingroup
		\centering
		ARTICLE \Roman{section}\par
		\vspace{0.3em}%
		\uline{\uppercase{#4}}\par
		\endgroup
	}{%
		#3#4%
	}%
}
\makeatother

\RedeclareSectionCommand[
	beforeskip=2ex plus 1ex minus .2ex,
	afterskip=3ex plus .2ex,
	font=\normalsize\bfseries
]{section}

% ---- Page styling ----------------------------------------------------------
\clearpairofpagestyles
\cfoot{-- \thepage{} of \ztotpages --}

% ---- Subsections (run-in, underlined, trailing period, dotted numbering) ----
\makeatletter
\renewcommand\subsection{\@startsection{subsection}{2}{\z@}%
	{-3.25ex\@plus -1ex \@minus -.2ex}%
	{-1em}% Negative value makes it run-in
	{\normalfont\normalsize\bfseries}}
\makeatother
\renewcommand{\thesubsection}{\arabic{section}.\arabic{subsection}}
\let\originalsubsection\subsection
\renewcommand{\subsection}[1]{%
	\originalsubsection{\uline{#1}.}%
}

% ---- Subsubsections (indented run-in) --------------------------------------
\makeatletter
\renewcommand\subsubsection{\@startsection{subsubsection}{3}{\parindent}%
	{-3.25ex\@plus -1ex \@minus -.2ex}%
	{-1em}%
	{\normalfont\normalsize\bfseries}}
\makeatother
\renewcommand{\thesubsubsection}{\arabic{section}.\arabic{subsection}.\arabic{subsubsection}}

\newlength{\subsubindent}
\setlength{\subsubindent}{\parindent}
\let\originalsubsubsection\subsubsection
\renewcommand{\subsubsection}[1]{%
	\originalsubsubsection{\uline{#1}.}%
	\parshape 2 \subsubindent \dimexpr\linewidth-\subsubindent\relax \subsubindent \dimexpr\linewidth-\subsubindent\relax%
}

% Hand-rolled 4th level.
\newcounter{subsubsubsection}[subsubsection]
\renewcommand{\thesubsubsubsection}{\arabic{section}.\arabic{subsection}.\arabic{subsubsection}.\arabic{subsubsubsection}}
\newcommand{\subsubsubsection}[1]{%
  \par\vspace{1ex}%
  \refstepcounter{subsubsubsection}%
  \noindent\hspace{2\parindent}%
  \textbf{\thesubsubsubsection~\uline{#1}.}~\ignorespaces%
  \parshape 2 2\subsubindent \dimexpr\linewidth-2\subsubindent\relax 2\subsubindent \dimexpr\linewidth-2\subsubindent\relax%
}

% ---- Table of Contents -----------------------------------------------------
\addto\captionsenglish{%
  \renewcommand{\contentsname}{Table of Contents}%
}
\renewcommand{\thesection}{\Roman{section}}
\DeclareTOCStyleEntry[
    entrynumberformat=\articleformat,
    numwidth=8em
]{tocline}{section}
\DeclareTOCStyleEntry[numwidth=3em, numsep=0.5em]{tocline}{subsection}
\DeclareTOCStyleEntry[numwidth=4em, numsep=0.5em]{tocline}{subsubsection}
\newcommand{\articleformat}[1]{Article #1}
\makeatletter
\renewcommand\tableofcontents{%
  \begingroup
  \setlength{\parindent}{0pt}
  \setlength{\parskip}{0pt}
  {\centering\Large\bfseries\contentsname\par}
  \vspace{1em}
  \@starttoc{toc}%
  \endgroup
}
\makeatother

% ---- hyperref loads LAST (clickable TOC / cross-refs / defined terms) -------
\usepackage[hidelinks]{hyperref}
  ->  \subimport{./src}{defined_terms}
%   ->  \subimport{./src}{variables}  ->  \subimport{./src}{content}
% Requires \documentclass{scrartcl} (KOMA-Script).
% ============================================================================

% ---- Compatible packages --------------------------------------------------
\usepackage{tocbasic}
\usepackage[english]{babel}
\usepackage{scrlayer-scrpage}
\usepackage{xcolor}
\usepackage{ulem}
\usepackage{hanging}
\usepackage{ifthen}
\usepackage[totpages]{zref}
\usepackage{tabularx}
\usepackage{array}
\usepackage{booktabs}
\usepackage{longtable}
\usepackage{xltabular}
\usepackage{enumitem}

% ---- Century Schoolbook (the typeface used by the U.S. Supreme Court) ------
\usepackage[T1]{fontenc}
\usepackage{tgschola}
\usepackage{xspace}

% Headings and title in the serif family (Century Schoolbook), not KOMA's sans.
\addtokomafont{disposition}{\rmfamily}

% ---- Document-wide constants ----------------------------------------------
% main.tex normally sets \title directly; \documentTitle is a safe fallback.
\providecommand{\documentTitle}{Document}
\newcommand{\omitText}{Intentionally Deleted}

% ---- Defined-term machinery -----------------------------------------------
% \dtregister stores a term's display text; \dtuse renders it in small caps and
% links back to its definition; \dtdef marks the definition (the anchor).
% The generated src/defined_terms.tex (\dtregister lines + \<key> uses) is
% \input by main.tex AFTER this file, so these commands exist when it loads.
\newcommand{\dtregister}[2]{\expandafter\def\csname dtdisplay@#1\endcsname{#2}}
\newcommand{\dtuse}[1]{\hyperlink{def:#1}{\textsc{\csname dtdisplay@#1\endcsname}}\xspace}
\newcommand{\dtdef}[1]{\hypertarget{def:#1}{\textsc{\csname dtdisplay@#1\endcsname}}\xspace}

% ---- Reusable helpers -----------------------------------------------------
% Centered bold non-section header (e.g. "RECITALS", "WITNESSETH").
\newcommand{\formaltext}[1]{%
	\begingroup
	\centering
	\normalfont\normalsize\bfseries #1\par
	\endgroup
	\vspace{0.5em}
}

\newcommand{\blankspace}{\ \newline \ }

% Inline fill-in blank (for the many "____" spots in these forms).
\newcommand{\fillin}[1][2.5cm]{\rule{#1}{0.4pt}}

% A signature rule (default 3in) with an optional caption line beneath it.
\newcommand{\sigline}[1][3in]{\rule{#1}{0.4pt}}
% \sigblock{caption} -> a signature rule with a caption typeset below it.
\newcommand{\sigblock}[1]{%
	\par\noindent\sigline\par\nobreak\vspace{-0.3em}\noindent #1\par
}

% Two-up signature row: a full-width rule with two captions beneath it (left/right).
% Shared by the board/shareholder consents, whose signature grid lists the directors
% or shareholders two per line above a common signature rule.
\newcommand{\sigpair}[2]{%
	\par\vspace{2em}\noindent\rule{\textwidth}{0.4pt}\par\nobreak\vspace{0.2em}%
	\noindent\makebox[0.5\textwidth][l]{#1}\makebox[0.5\textwidth][l]{#2}\par
}

% Exhibit / Schedule heading: centered "Exhibit <letter>" + underlined title.
% Used instead of \section so exhibits stay out of the article numbering and TOC.
\newcommand{\exhibithead}[2]{%
	\begingroup
	\centering
	\normalfont\large\bfseries Exhibit #1\par
	\vspace{0.3em}
	\uline{\uppercase{#2}}\par
	\endgroup
	\vspace{1em}
}
\newcommand{\schedulehead}[2]{%
	\begingroup
	\centering
	\normalfont\large\bfseries Schedule #1\par
	\vspace{0.3em}
	\uline{\uppercase{#2}}\par
	\endgroup
	\vspace{1em}
}

% ---- Title -----------------------------------------------------------------
% Completely redefine \maketitle to avoid any KOMA errors and center a bold title.
\makeatletter
\renewcommand{\maketitle}{%
	\begingroup
	\centering
	\normalfont\large\bfseries\color{black!80}%
	\@title\par
	\vspace{0.5em}
	\endgroup
}
\makeatother

% ---- Sections: centered "ARTICLE <Roman>" + underlined uppercase title ------
\makeatletter
\renewcommand{\sectionlinesformat}[4]{%
	\Ifstr{#1}{section}{%
		\begingroup
		\centering
		ARTICLE \Roman{section}\par
		\vspace{0.3em}%
		\uline{\uppercase{#4}}\par
		\endgroup
	}{%
		#3#4%
	}%
}
\makeatother

\RedeclareSectionCommand[
	beforeskip=2ex plus 1ex minus .2ex,
	afterskip=3ex plus .2ex,
	font=\normalsize\bfseries
]{section}

% ---- Page styling ----------------------------------------------------------
\clearpairofpagestyles
\cfoot{-- \thepage{} of \ztotpages --}

% ---- Subsections (run-in, underlined, trailing period, dotted numbering) ----
\makeatletter
\renewcommand\subsection{\@startsection{subsection}{2}{\z@}%
	{-3.25ex\@plus -1ex \@minus -.2ex}%
	{-1em}% Negative value makes it run-in
	{\normalfont\normalsize\bfseries}}
\makeatother
\renewcommand{\thesubsection}{\arabic{section}.\arabic{subsection}}
\let\originalsubsection\subsection
\renewcommand{\subsection}[1]{%
	\originalsubsection{\uline{#1}.}%
}

% ---- Subsubsections (indented run-in) --------------------------------------
\makeatletter
\renewcommand\subsubsection{\@startsection{subsubsection}{3}{\parindent}%
	{-3.25ex\@plus -1ex \@minus -.2ex}%
	{-1em}%
	{\normalfont\normalsize\bfseries}}
\makeatother
\renewcommand{\thesubsubsection}{\arabic{section}.\arabic{subsection}.\arabic{subsubsection}}

\newlength{\subsubindent}
\setlength{\subsubindent}{\parindent}
\let\originalsubsubsection\subsubsection
\renewcommand{\subsubsection}[1]{%
	\originalsubsubsection{\uline{#1}.}%
	\parshape 2 \subsubindent \dimexpr\linewidth-\subsubindent\relax \subsubindent \dimexpr\linewidth-\subsubindent\relax%
}

% Hand-rolled 4th level.
\newcounter{subsubsubsection}[subsubsection]
\renewcommand{\thesubsubsubsection}{\arabic{section}.\arabic{subsection}.\arabic{subsubsection}.\arabic{subsubsubsection}}
\newcommand{\subsubsubsection}[1]{%
  \par\vspace{1ex}%
  \refstepcounter{subsubsubsection}%
  \noindent\hspace{2\parindent}%
  \textbf{\thesubsubsubsection~\uline{#1}.}~\ignorespaces%
  \parshape 2 2\subsubindent \dimexpr\linewidth-2\subsubindent\relax 2\subsubindent \dimexpr\linewidth-2\subsubindent\relax%
}

% ---- Table of Contents -----------------------------------------------------
\addto\captionsenglish{%
  \renewcommand{\contentsname}{Table of Contents}%
}
\renewcommand{\thesection}{\Roman{section}}
\DeclareTOCStyleEntry[
    entrynumberformat=\articleformat,
    numwidth=8em
]{tocline}{section}
\DeclareTOCStyleEntry[numwidth=3em, numsep=0.5em]{tocline}{subsection}
\DeclareTOCStyleEntry[numwidth=4em, numsep=0.5em]{tocline}{subsubsection}
\newcommand{\articleformat}[1]{Article #1}
\makeatletter
\renewcommand\tableofcontents{%
  \begingroup
  \setlength{\parindent}{0pt}
  \setlength{\parskip}{0pt}
  {\centering\Large\bfseries\contentsname\par}
  \vspace{1em}
  \@starttoc{toc}%
  \endgroup
}
\makeatother

% ---- hyperref loads LAST (clickable TOC / cross-refs / defined terms) -------
\usepackage[hidelinks]{hyperref}
  ->  \subimport{./src}{defined_terms}
%   ->  \subimport{./src}{variables}  ->  \subimport{./src}{content}
% Requires \documentclass{scrartcl} (KOMA-Script).
% ============================================================================

% ---- Compatible packages --------------------------------------------------
\usepackage{tocbasic}
\usepackage[english]{babel}
\usepackage{scrlayer-scrpage}
\usepackage{xcolor}
\usepackage{ulem}
\usepackage{hanging}
\usepackage{ifthen}
\usepackage[totpages]{zref}
\usepackage{tabularx}
\usepackage{array}
\usepackage{booktabs}
\usepackage{longtable}
\usepackage{xltabular}
\usepackage{enumitem}

% ---- Century Schoolbook (the typeface used by the U.S. Supreme Court) ------
\usepackage[T1]{fontenc}
\usepackage{tgschola}
\usepackage{xspace}

% Headings and title in the serif family (Century Schoolbook), not KOMA's sans.
\addtokomafont{disposition}{\rmfamily}

% ---- Document-wide constants ----------------------------------------------
% main.tex normally sets \title directly; \documentTitle is a safe fallback.
\providecommand{\documentTitle}{Document}
\newcommand{\omitText}{Intentionally Deleted}

% ---- Defined-term machinery -----------------------------------------------
% \dtregister stores a term's display text; \dtuse renders it in small caps and
% links back to its definition; \dtdef marks the definition (the anchor).
% The generated src/defined_terms.tex (\dtregister lines + \<key> uses) is
% \input by main.tex AFTER this file, so these commands exist when it loads.
\newcommand{\dtregister}[2]{\expandafter\def\csname dtdisplay@#1\endcsname{#2}}
\newcommand{\dtuse}[1]{\hyperlink{def:#1}{\textsc{\csname dtdisplay@#1\endcsname}}\xspace}
\newcommand{\dtdef}[1]{\hypertarget{def:#1}{\textsc{\csname dtdisplay@#1\endcsname}}\xspace}

% ---- Reusable helpers -----------------------------------------------------
% Centered bold non-section header (e.g. "RECITALS", "WITNESSETH").
\newcommand{\formaltext}[1]{%
	\begingroup
	\centering
	\normalfont\normalsize\bfseries #1\par
	\endgroup
	\vspace{0.5em}
}

\newcommand{\blankspace}{\ \newline \ }

% Inline fill-in blank (for the many "____" spots in these forms).
\newcommand{\fillin}[1][2.5cm]{\rule{#1}{0.4pt}}

% A signature rule (default 3in) with an optional caption line beneath it.
\newcommand{\sigline}[1][3in]{\rule{#1}{0.4pt}}
% \sigblock{caption} -> a signature rule with a caption typeset below it.
\newcommand{\sigblock}[1]{%
	\par\noindent\sigline\par\nobreak\vspace{-0.3em}\noindent #1\par
}

% Two-up signature row: a full-width rule with two captions beneath it (left/right).
% Shared by the board/shareholder consents, whose signature grid lists the directors
% or shareholders two per line above a common signature rule.
\newcommand{\sigpair}[2]{%
	\par\vspace{2em}\noindent\rule{\textwidth}{0.4pt}\par\nobreak\vspace{0.2em}%
	\noindent\makebox[0.5\textwidth][l]{#1}\makebox[0.5\textwidth][l]{#2}\par
}

% Exhibit / Schedule heading: centered "Exhibit <letter>" + underlined title.
% Used instead of \section so exhibits stay out of the article numbering and TOC.
\newcommand{\exhibithead}[2]{%
	\begingroup
	\centering
	\normalfont\large\bfseries Exhibit #1\par
	\vspace{0.3em}
	\uline{\uppercase{#2}}\par
	\endgroup
	\vspace{1em}
}
\newcommand{\schedulehead}[2]{%
	\begingroup
	\centering
	\normalfont\large\bfseries Schedule #1\par
	\vspace{0.3em}
	\uline{\uppercase{#2}}\par
	\endgroup
	\vspace{1em}
}

% ---- Title -----------------------------------------------------------------
% Completely redefine \maketitle to avoid any KOMA errors and center a bold title.
\makeatletter
\renewcommand{\maketitle}{%
	\begingroup
	\centering
	\normalfont\large\bfseries\color{black!80}%
	\@title\par
	\vspace{0.5em}
	\endgroup
}
\makeatother

% ---- Sections: centered "ARTICLE <Roman>" + underlined uppercase title ------
\makeatletter
\renewcommand{\sectionlinesformat}[4]{%
	\Ifstr{#1}{section}{%
		\begingroup
		\centering
		ARTICLE \Roman{section}\par
		\vspace{0.3em}%
		\uline{\uppercase{#4}}\par
		\endgroup
	}{%
		#3#4%
	}%
}
\makeatother

\RedeclareSectionCommand[
	beforeskip=2ex plus 1ex minus .2ex,
	afterskip=3ex plus .2ex,
	font=\normalsize\bfseries
]{section}

% ---- Page styling ----------------------------------------------------------
\clearpairofpagestyles
\cfoot{-- \thepage{} of \ztotpages --}

% ---- Subsections (run-in, underlined, trailing period, dotted numbering) ----
\makeatletter
\renewcommand\subsection{\@startsection{subsection}{2}{\z@}%
	{-3.25ex\@plus -1ex \@minus -.2ex}%
	{-1em}% Negative value makes it run-in
	{\normalfont\normalsize\bfseries}}
\makeatother
\renewcommand{\thesubsection}{\arabic{section}.\arabic{subsection}}
\let\originalsubsection\subsection
\renewcommand{\subsection}[1]{%
	\originalsubsection{\uline{#1}.}%
}

% ---- Subsubsections (indented run-in) --------------------------------------
\makeatletter
\renewcommand\subsubsection{\@startsection{subsubsection}{3}{\parindent}%
	{-3.25ex\@plus -1ex \@minus -.2ex}%
	{-1em}%
	{\normalfont\normalsize\bfseries}}
\makeatother
\renewcommand{\thesubsubsection}{\arabic{section}.\arabic{subsection}.\arabic{subsubsection}}

\newlength{\subsubindent}
\setlength{\subsubindent}{\parindent}
\let\originalsubsubsection\subsubsection
\renewcommand{\subsubsection}[1]{%
	\originalsubsubsection{\uline{#1}.}%
	\parshape 2 \subsubindent \dimexpr\linewidth-\subsubindent\relax \subsubindent \dimexpr\linewidth-\subsubindent\relax%
}

% Hand-rolled 4th level.
\newcounter{subsubsubsection}[subsubsection]
\renewcommand{\thesubsubsubsection}{\arabic{section}.\arabic{subsection}.\arabic{subsubsection}.\arabic{subsubsubsection}}
\newcommand{\subsubsubsection}[1]{%
  \par\vspace{1ex}%
  \refstepcounter{subsubsubsection}%
  \noindent\hspace{2\parindent}%
  \textbf{\thesubsubsubsection~\uline{#1}.}~\ignorespaces%
  \parshape 2 2\subsubindent \dimexpr\linewidth-2\subsubindent\relax 2\subsubindent \dimexpr\linewidth-2\subsubindent\relax%
}

% ---- Table of Contents -----------------------------------------------------
\addto\captionsenglish{%
  \renewcommand{\contentsname}{Table of Contents}%
}
\renewcommand{\thesection}{\Roman{section}}
\DeclareTOCStyleEntry[
    entrynumberformat=\articleformat,
    numwidth=8em
]{tocline}{section}
\DeclareTOCStyleEntry[numwidth=3em, numsep=0.5em]{tocline}{subsection}
\DeclareTOCStyleEntry[numwidth=4em, numsep=0.5em]{tocline}{subsubsection}
\newcommand{\articleformat}[1]{Article #1}
\makeatletter
\renewcommand\tableofcontents{%
  \begingroup
  \setlength{\parindent}{0pt}
  \setlength{\parskip}{0pt}
  {\centering\Large\bfseries\contentsname\par}
  \vspace{1em}
  \@starttoc{toc}%
  \endgroup
}
\makeatother

% ---- hyperref loads LAST (clickable TOC / cross-refs / defined terms) -------
\usepackage[hidelinks]{hyperref}
  ->  \subimport{./src}{defined_terms}
%   ->  \subimport{./src}{variables}  ->  \subimport{./src}{content}
% Requires \documentclass{scrartcl} (KOMA-Script).
% ============================================================================

% ---- Compatible packages --------------------------------------------------
\usepackage{tocbasic}
\usepackage[english]{babel}
\usepackage{scrlayer-scrpage}
\usepackage{xcolor}
\usepackage{ulem}
\usepackage{hanging}
\usepackage{ifthen}
\usepackage[totpages]{zref}
\usepackage{tabularx}
\usepackage{array}
\usepackage{booktabs}
\usepackage{longtable}
\usepackage{xltabular}
\usepackage{enumitem}

% ---- Century Schoolbook (the typeface used by the U.S. Supreme Court) ------
\usepackage[T1]{fontenc}
\usepackage{tgschola}
\usepackage{xspace}

% Headings and title in the serif family (Century Schoolbook), not KOMA's sans.
\addtokomafont{disposition}{\rmfamily}

% ---- Document-wide constants ----------------------------------------------
% main.tex normally sets \title directly; \documentTitle is a safe fallback.
\providecommand{\documentTitle}{Document}
\newcommand{\omitText}{Intentionally Deleted}

% ---- Defined-term machinery -----------------------------------------------
% \dtregister stores a term's display text; \dtuse renders it in small caps and
% links back to its definition; \dtdef marks the definition (the anchor).
% The generated src/defined_terms.tex (\dtregister lines + \<key> uses) is
% \input by main.tex AFTER this file, so these commands exist when it loads.
\newcommand{\dtregister}[2]{\expandafter\def\csname dtdisplay@#1\endcsname{#2}}
\newcommand{\dtuse}[1]{\hyperlink{def:#1}{\textsc{\csname dtdisplay@#1\endcsname}}\xspace}
\newcommand{\dtdef}[1]{\hypertarget{def:#1}{\textsc{\csname dtdisplay@#1\endcsname}}\xspace}

% ---- Reusable helpers -----------------------------------------------------
% Centered bold non-section header (e.g. "RECITALS", "WITNESSETH").
\newcommand{\formaltext}[1]{%
	\begingroup
	\centering
	\normalfont\normalsize\bfseries #1\par
	\endgroup
	\vspace{0.5em}
}

\newcommand{\blankspace}{\ \newline \ }

% Inline fill-in blank (for the many "____" spots in these forms).
\newcommand{\fillin}[1][2.5cm]{\rule{#1}{0.4pt}}

% A signature rule (default 3in) with an optional caption line beneath it.
\newcommand{\sigline}[1][3in]{\rule{#1}{0.4pt}}
% \sigblock{caption} -> a signature rule with a caption typeset below it.
\newcommand{\sigblock}[1]{%
	\par\noindent\sigline\par\nobreak\vspace{-0.3em}\noindent #1\par
}

% Two-up signature row: a full-width rule with two captions beneath it (left/right).
% Shared by the board/shareholder consents, whose signature grid lists the directors
% or shareholders two per line above a common signature rule.
\newcommand{\sigpair}[2]{%
	\par\vspace{2em}\noindent\rule{\textwidth}{0.4pt}\par\nobreak\vspace{0.2em}%
	\noindent\makebox[0.5\textwidth][l]{#1}\makebox[0.5\textwidth][l]{#2}\par
}

% Exhibit / Schedule heading: centered "Exhibit <letter>" + underlined title.
% Used instead of \section so exhibits stay out of the article numbering and TOC.
\newcommand{\exhibithead}[2]{%
	\begingroup
	\centering
	\normalfont\large\bfseries Exhibit #1\par
	\vspace{0.3em}
	\uline{\uppercase{#2}}\par
	\endgroup
	\vspace{1em}
}
\newcommand{\schedulehead}[2]{%
	\begingroup
	\centering
	\normalfont\large\bfseries Schedule #1\par
	\vspace{0.3em}
	\uline{\uppercase{#2}}\par
	\endgroup
	\vspace{1em}
}

% ---- Title -----------------------------------------------------------------
% Completely redefine \maketitle to avoid any KOMA errors and center a bold title.
\makeatletter
\renewcommand{\maketitle}{%
	\begingroup
	\centering
	\normalfont\large\bfseries\color{black!80}%
	\@title\par
	\vspace{0.5em}
	\endgroup
}
\makeatother

% ---- Sections: centered "ARTICLE <Roman>" + underlined uppercase title ------
\makeatletter
\renewcommand{\sectionlinesformat}[4]{%
	\Ifstr{#1}{section}{%
		\begingroup
		\centering
		ARTICLE \Roman{section}\par
		\vspace{0.3em}%
		\uline{\uppercase{#4}}\par
		\endgroup
	}{%
		#3#4%
	}%
}
\makeatother

\RedeclareSectionCommand[
	beforeskip=2ex plus 1ex minus .2ex,
	afterskip=3ex plus .2ex,
	font=\normalsize\bfseries
]{section}

% ---- Page styling ----------------------------------------------------------
\clearpairofpagestyles
\cfoot{-- \thepage{} of \ztotpages --}

% ---- Subsections (run-in, underlined, trailing period, dotted numbering) ----
\makeatletter
\renewcommand\subsection{\@startsection{subsection}{2}{\z@}%
	{-3.25ex\@plus -1ex \@minus -.2ex}%
	{-1em}% Negative value makes it run-in
	{\normalfont\normalsize\bfseries}}
\makeatother
\renewcommand{\thesubsection}{\arabic{section}.\arabic{subsection}}
\let\originalsubsection\subsection
\renewcommand{\subsection}[1]{%
	\originalsubsection{\uline{#1}.}%
}

% ---- Subsubsections (indented run-in) --------------------------------------
\makeatletter
\renewcommand\subsubsection{\@startsection{subsubsection}{3}{\parindent}%
	{-3.25ex\@plus -1ex \@minus -.2ex}%
	{-1em}%
	{\normalfont\normalsize\bfseries}}
\makeatother
\renewcommand{\thesubsubsection}{\arabic{section}.\arabic{subsection}.\arabic{subsubsection}}

\newlength{\subsubindent}
\setlength{\subsubindent}{\parindent}
\let\originalsubsubsection\subsubsection
\renewcommand{\subsubsection}[1]{%
	\originalsubsubsection{\uline{#1}.}%
	\parshape 2 \subsubindent \dimexpr\linewidth-\subsubindent\relax \subsubindent \dimexpr\linewidth-\subsubindent\relax%
}

% Hand-rolled 4th level.
\newcounter{subsubsubsection}[subsubsection]
\renewcommand{\thesubsubsubsection}{\arabic{section}.\arabic{subsection}.\arabic{subsubsection}.\arabic{subsubsubsection}}
\newcommand{\subsubsubsection}[1]{%
  \par\vspace{1ex}%
  \refstepcounter{subsubsubsection}%
  \noindent\hspace{2\parindent}%
  \textbf{\thesubsubsubsection~\uline{#1}.}~\ignorespaces%
  \parshape 2 2\subsubindent \dimexpr\linewidth-2\subsubindent\relax 2\subsubindent \dimexpr\linewidth-2\subsubindent\relax%
}

% ---- Table of Contents -----------------------------------------------------
\addto\captionsenglish{%
  \renewcommand{\contentsname}{Table of Contents}%
}
\renewcommand{\thesection}{\Roman{section}}
\DeclareTOCStyleEntry[
    entrynumberformat=\articleformat,
    numwidth=8em
]{tocline}{section}
\DeclareTOCStyleEntry[numwidth=3em, numsep=0.5em]{tocline}{subsection}
\DeclareTOCStyleEntry[numwidth=4em, numsep=0.5em]{tocline}{subsubsection}
\newcommand{\articleformat}[1]{Article #1}
\makeatletter
\renewcommand\tableofcontents{%
  \begingroup
  \setlength{\parindent}{0pt}
  \setlength{\parskip}{0pt}
  {\centering\Large\bfseries\contentsname\par}
  \vspace{1em}
  \@starttoc{toc}%
  \endgroup
}
\makeatother

% ---- hyperref loads LAST (clickable TOC / cross-refs / defined terms) -------
\usepackage[hidelinks]{hyperref}
